\section{Compression by cycles and observations}\label{sec:compression}
This section compresses the exact block formulation in two steps. First,
cycle coordinates replace arc coordinates and suppress long paths. Second,
observations determine some cycle coordinates, leaving only cycles supported
on unobserved arcs. We assume $\Flow\ne\varnothing$; if the original flow
domain is empty, its hull is described by any inconsistent linear system.

\subsection{Fixed coordinates and reduced graphs}\label{subsec:fixed-arcs}
The structural parameters can be applied after elementary exact preprocessing.
This is useful when capacities force a nominal cycle to carry no freedom.

\begin{lemma}[Deleting fixed arcs]\label{lem:fixed-arcs}
Let $F\subseteq E$ and suppose $x_e=c_e$ throughout $\Flow$ for every
$e\in F$, where $c\in\Q^F$ is known. Set $R=E\setminus F$,
$b'=b-A_{:,F}c$, and
\[
 \Flow'=\{x_R:A_{:,R}x_R=b',\ 0\le x_R\le u_R\},\qquad
 \Obs'=\Obs\cap(R\times[m]).
\]
The hull for $(\Flow',\Obs')$, together with
$x_F=c$ and $z_{ej}=c_e y_j$ for observed $e\in F$, is affinely
isomorphic to $\Hull$. If $F$ contains every coordinate fixed throughout
$\Flow$, then
\begin{equation}\label{eq:reduced-affine-hull}
 \aff\Flow'=\{x_R:A_{:,R}x_R=b'\}.
\end{equation}
\end{lemma}
\begin{proof}
Feasible vectors of $\Flow$ restrict to $\Flow'$. Conversely, adjoining
$c$ to a vector of $\Flow'$ satisfies the original balances and bounds,
because $0\le c\le u_F$. The same bijection at graph points adjoins the
affine coordinates $z_{ej}=c_e y_j$, so it commutes with convexification.

For the last assertion, every remaining coordinate is nonconstant on
$\Flow'$. For each $e\in R$, choose two feasible vectors differing in
coordinate $e$; their midpoint has $0<x_e<u_e$. Average these midpoints
over $e\in R$. The resulting feasible vector is strictly inside every
remaining bound. A sufficiently small relative neighborhood of that vector
in the balance affine space therefore remains feasible, proving
\eqref{eq:reduced-affine-hull}. If $R$ is empty, both spaces are the same
singleton and the assertion is immediate.
\end{proof}

The lemma takes the fixed coordinates as input; it does not require a
particular detection algorithm. Deleting them can split blocks or reduce
their cycle ranks. Once all fixed coordinates are removed, the remaining
cycle space is the direction space of the feasible affine hull. Conditional
slices at a particular observation or zero state weight can still have
smaller dimension.

\subsection{Suppressing paths}\label{subsec:suppression}
Fix a cyclic block $B$. Degrees in this subsection are degrees \emph{within
the block}, with loops counted twice. For $r_B\ge2$, suppress maximal paths
whose internal vertices have degree two, and give each resulting core edge
an arbitrary orientation. A rank-one block is a single cycle; represent
its entire cyclic traversal by one core loop. This includes two parallel
edges and an original self-loop.

Write $\mathcal P_B$ for the core edges, or equivalently the original
paths, and $k_B=|\mathcal P_B|$. For an original arc $e$ on path $p$, let
$\varepsilon_e\in\{-1,1\}$ compare its orientation with the chosen
path traversal. A circulation has constant signed deviation on a path:
\begin{equation}\label{eq:path-deviation}
 x_e=v_e+\varepsilon_e\delta_p\qquad(e\in p).
\end{equation}
The reason is balance of the block circulation at each internal degree-two
vertex. Balances at the unsuppressed vertices say precisely that $\delta$
is a circulation of the core. Arc bounds give the path interval
\begin{align}
 L_p&=\max_{e\in p}\ell_e,& U_p&=\min_{e\in p}u'_e,\label{eq:path-intervals}\\
 [\ell_e,u'_e]&=
 \begin{cases}[-v_e,u_e-v_e],&\varepsilon_e=1,\\
 [v_e-u_e,v_e],&\varepsilon_e=-1.
 \end{cases}\nonumber
\end{align}
All intervals are finite. Their joint feasibility follows from
$\Flow\ne\varnothing$.

Suppression preserves cycle rank. If $r_B\ge2$, the connected core has
minimum degree at least three: otherwise the original block would have
been a single cycle. If its vertex count is $n_B$, then
$2k_B\ge3n_B$ and $r_B=k_B-n_B+1$, so
\begin{equation}\label{eq:core-size}
 n_B\le2r_B-2,\qquad k_B\le3r_B-3.
\end{equation}
For rank one, $k_B=1$. In particular, $k_B\le3r_B$ in every case.

Choose a spanning tree of the core and let
$\widehat C_B\in\{0,\pm1\}^{k_B\times r_B}$ be its fundamental-cycle
matrix, normalized to have identity chord rows. Its columns parametrize
core circulations bijectively. Thus the feasible core domain is
\begin{equation}\label{eq:core-domain}
 T_B=\{t\in\R^{r_B}:L_p\le(\widehat C_Bt)_p\le U_p
                         \quad(p\in\mathcal P_B)\}.
\end{equation}
It is nonempty and bounded, including in lower-dimensional cases. Boundedness
follows already from the chord rows. For later use, choose one original
representative arc $e(p)$ on each path and put
\begin{equation}\label{eq:aggregate-path-expression}
 s_p(x)=\varepsilon_{e(p)}(x_{e(p)}-v_{e(p)}),\qquad
 q_{ej}=\varepsilon_e(z_{ej}-v_e y_j).
\end{equation}
For $x\in\Flow$, $s_p(x)$ is its core deviation. A separate observation
$q_{ej}$ is retained for \emph{every} observed original arc, including
multiple observations of one path.

\begin{theorem}[Cycle-coordinate extended hull]\label{thm:cycle-compression}
The hull $\Hull$ has an exact rational extended formulation with
\begin{equation}\label{eq:initial-variable-count}
 \sum_{B\in\Blocks}r_Ba_B
\end{equation}
additional variables and
\begin{equation}\label{eq:compression-row-count}
 O\left(|V|+|E|+m+|\Obs|
       +\sum_{B:a_B>0}r_B(a_B+1)\right)
\end{equation}
rows, counting the original domain constraints. Coefficients of $x$, $z$,
and additional variables can all be chosen from $\{0,\pm1\}$. Coefficients
of $y$ and the right-hand sides may depend on the rational network data.
All coefficient encoding lengths are polynomial in the input length.
\end{theorem}
\begin{proof}
Retain $x\in\Flow$, $y\in\Delta_m$, and the bridge equations
\eqref{eq:bridge-product}. For each $B$ with $a_B>0$, introduce
$t^{B,j}\in\R^{r_B}$ for $j\in J_B$ and impose
\begin{align}
 y_jL_p&\le(\widehat C_Bt^{B,j})_p\le y_jU_p,
       &&j\in J_B,\ p\in\mathcal P_B,\label{eq:compressed-state}\\
 q_{ej}&=(\widehat C_Bt^{B,j})_p,
       &&(e,j)\in\Obs,\ e\in p\subseteq E_B,\label{eq:compressed-observation}\\
 \lambda_{B,*}L_p&\le s_p(x)-\sum_{j\in J_B}
                      (\widehat C_Bt^{B,j})_p
                    \le\lambda_{B,*}U_p,
       &&p\in\mathcal P_B.\label{eq:compressed-residual}
\end{align}
The aggregate deviation $s(x)$ is a core circulation. Hence the middle
quantity in \eqref{eq:compressed-residual} is a core circulation as well.
These constraints are exactly \cref{thm:block-state-reduction} in core
coordinates, with the merged state eliminated by its sum equation.
At a zero weight, the finite path bounds force every core path value to
zero; the identity chord rows force the entire state vector to zero.
Consequently the equivalence includes all simplex boundary strata. Blocks
with $a_B=0$ add no constraints, as in the block theorem.

There are $r_Ba_B$ variables, $2k_B(a_B+1)$ bound rows, and one equation
per block observation. Bound \eqref{eq:core-size} proves the row count.
Every state variable has coefficient $0,1$, or $-1$. An observation row
has one original product coordinate, and each residual row uses one
original flow coordinate through $s_p(x)$; no sums of different cycle
expressions for that flow coordinate are necessary. This proves the stated
coefficient bound. Reference flows are sums of rational balances; path
endpoints are selected signed capacity offsets. The remaining coefficients
are sums of these data, with polynomial rational encoding length.
\end{proof}

\subsection{Eliminating observed directions}\label{subsec:observed-directions}
Define the observed and unobserved arcs for a block and label by
\[
 O_{Bj}=\{e\in E_B:(e,j)\in\Obs\},\qquad U_{Bj}=E_B\setminus O_{Bj}.
\]
Let $c(V_B,U_{Bj})$ count all connected components of this unobserved
undirected subgraph, including isolated vertices. Its cycle rank is
\begin{equation}\label{eq:unobserved-rank}
 \rho_{Bj}=|U_{Bj}|-|V_B|+c(V_B,U_{Bj}).
\end{equation}
Loops and parallel edges keep their usual multigraph meanings. In particular,
each unobserved loop contributes one to the cycle rank.

\begin{lemma}[Unit minors for cycle elimination]\label{lem:cycle-tu}
A fundamental-cycle matrix with identity chord rows is totally unimodular.
Let $C$ be such a matrix, let $D$ consist of $d$ linearly independent rows,
and choose $d$ columns $I$ such that $B=D_{:,I}$ is nonsingular. Write $F$
for the complementary columns. Then, for every row $c$ of $C$, the vectors
\begin{equation}\label{eq:unit-schur}
 W=c_I B^{-1},\qquad R=c_F-WD_{:,F}
\end{equation}
have entries in $\{0,\pm1\}$.
\end{lemma}
\begin{proof}
Delete one incidence row of the connected core and partition the reduced
incidence matrix as $[N_T\ N_Q]$ into tree and chord columns. A tree
basis has determinant $\pm1$. Every square incidence minor is $0$ or
$\pm1$: induction expands a column with at most one nonzero entry; if
every column has two entries, they are opposite and the row sum is zero.
Every minor of $N_T^{-1}N_Q$ equals, up to sign, the determinant of a
matrix obtained by replacing the corresponding tree columns by chord
columns, divided by $\det N_T$. This follows by multiplying the replacement
matrix by $N_T^{-1}$ and expanding its unchanged identity columns.
The replacement is a square incidence submatrix. Therefore
$-N_T^{-1}N_Q$ is totally unimodular. Appending identity chord rows
preserves this property, proving the assertion for $C$. A core consisting
of one loop has $C=(1)$ directly.

In particular $\det B=\pm1$. Cramer's rule applied to $WB=c_I$ expresses
each entry of $W$ as a selected-row replacement minor of $C$ divided by
$\det B$. For a free column $f$, the bordered determinant identity gives
\[
 R_f=\frac{\det\begin{pmatrix}B&D_{:,f}\\c_I&c_f\end{pmatrix}}
                  {\det B}.
\]
The numerator is a minor of $C$, or zero when its rows repeat. These
formulas prove \eqref{eq:unit-schur}. Empty pivot or free sets have their
usual zero-dimensional matrix interpretation.
\end{proof}

\begin{theorem}[Observation-sensitive extended hull]\label{thm:observed-compression}
There is an exact rational extended formulation for $\Hull$ with
\begin{equation}\label{eq:observed-variable-count}
 \sum_{B\in\Blocks}\ \sum_{j\in J_B}\rho_{Bj}
\end{equation}
additional variables and the row bound \eqref{eq:compression-row-count}.
Coefficients of $x$, $z$, and additional variables belong to $\{0,\pm1\}$;
simplex coefficients and right-hand sides have polynomial rational encoding
length. The variable count is an upper bound for this construction, not a
lower bound for arbitrary extended formulations.
\end{theorem}
\begin{proof}
Fix $B$ and $j\in J_B$, suppress their subscripts, and write
$C=\widehat C_B$, $r=r_B$. Choose a maximal independent set of observed
path rows of $C$, giving $D\in\R^{d\times r}$. For each selected row
choose one actual observed arc and use its value from
\eqref{eq:aggregate-path-expression}, giving $q\in\R^d$. The selected
entries of $q$ involve distinct original product variables. All other
observations will remain as equations, including any repeated path values.

Choose pivot columns $I$ and complementary columns $F$ as in
\cref{lem:cycle-tu}. Retain $g=t_F\in\R^{r-d}$ and eliminate the pivot
coordinates by the reversible substitution
\begin{equation}\label{eq:pivot-reconstruction}
 t_I=B^{-1}(q-D_{:,F}g).
\end{equation}
For every core path $p$ with row $c_p$ define
\begin{equation}\label{eq:path-reconstruction}
 \phi_{pj}=(c_{p,I}B^{-1})q+
            (c_{p,F}-c_{p,I}B^{-1}D_{:,F})g.
\end{equation}
Replace $(\widehat C_Bt^{B,j})_p$ by $\phi_{pj}$ in
\eqref{eq:compressed-state}--\eqref{eq:compressed-residual}. Selected
observation equations become identities and can be omitted; all other
observation equations remain. Formula \eqref{eq:pivot-reconstruction}
proves equivalence in both directions, so no projection condition is lost.

By \cref{lem:cycle-tu}, each selected product and retained variable has
coefficient $0,1$, or $-1$ in every $\phi_{pj}$. A nonselected observation
adds a distinct product coordinate; it cannot add a second occurrence of
a selected product. Different labels use disjoint product and auxiliary
columns. The residual uses one original flow coordinate through $s_p(x)$.
Thus collecting terms does not increase any $x$, $z$, or auxiliary
coefficient beyond one in magnitude. Only coefficients of $y$ collect
reference-flow and interval data. Their encoding lengths remain polynomial.
Elimination adds no rows.

It remains to identify $r-d$. Restrict a block circulation to $O_{Bj}$.
The kernel is precisely the space of circulations supported on $U_{Bj}$,
whose dimension is \eqref{eq:unobserved-rank} by the incidence-rank formula.
The core matrix parametrizes all block circulations, and observing any
arc of a suppressed path fixes its signed path value. Therefore the rank
of the restriction map is exactly the rank $d$ of the observed core rows.
Rank--nullity gives $r_B-d=\rho_{Bj}$, proving
\eqref{eq:observed-variable-count}.
\end{proof}

For clarity, the final block inequalities are
\begin{align}
 y_jL_p&\le\phi_{pj}\le y_jU_p,
       &&j\in J_B,\ p\in\mathcal P_B,\label{eq:observed-state-bounds}\\
 q_{ej}&=\phi_{pj},
       &&(e,j)\in\Obs,\ e\in p\subseteq E_B,\label{eq:observed-equalities}\\
 \lambda_{B,*}L_p&\le s_p(x)-\sum_{j\in J_B}\phi_{pj}
                       \le\lambda_{B,*}U_p,
       &&p\in\mathcal P_B.\label{eq:observed-residual-bounds}
\end{align}
The original flow and simplex domains and bridge observations remain part
of the formulation. In particular, reconstructing state products by balance
equations alone is insufficient: all displayed state and residual bounds
must also be imposed.

The row counts do not bound matrix nonzeros when ranks grow. If
$N=|V|+|E|+m+|\Obs|$ and $\Obs_B=\Obs\cap(E_B\times[m])$, both
constructions have the explicit nonzero bound
\begin{equation}\label{eq:compression-nonzeros}
 O\left(N+\sum_{B:a_B>0}
       \bigl[r_B^2(a_B+1)+r_B|\Obs_B|\bigr]\right).
\end{equation}
Indeed, a reconstructed path has $O(r_B)$ selected-product and auxiliary
terms per label; there are $O(r_B)$ paths and $a_B$ explicit labels.
Each remaining observation also uses $O(r_B)$ terms. Hence bounded maximum
block rank gives size linear in the sparse input, even as $m$ and the graph
grow. The whole formulation need not be totally unimodular.
The unit coefficient statements use the displayed rational row scaling;
clearing denominators in the simplex coefficients can destroy that unit
bound in a primitive integer normalization.

\subsection{Forest complements and product completion}\label{subsec:forest-complements}
\begin{corollary}[An original-coordinate hull on arbitrary graphs]\label{cor:forest-hull}
Suppose $(V_B,U_{Bj})$ is a forest for every $B$ and $j\in J_B$.
Then \eqref{eq:observed-state-bounds}--\eqref{eq:observed-residual-bounds},
the original domains, and the bridge equations describe $\Hull$ with no
additional variables. They have unit flow and product coefficients and
the row bound \eqref{eq:compression-row-count}.
\end{corollary}
\begin{proof}
An undirected multigraph is a forest exactly when its cycle rank is zero.
Thus every $\rho_{Bj}$ in \eqref{eq:observed-variable-count} vanishes.
Apply \cref{thm:observed-compression}.
\end{proof}

Equivalently, the observations for each locally observed label meet every
cycle of its block. Entirely unobserved labels impose no condition: they
are already included in the merged state. The condition is sufficient;
it does not assert that remaining unobserved cycles prevent an
original-coordinate description with unit coefficients.

\begin{proposition}[Minimum individual-coordinate completion]\label{prop:minimum-completion}
Fix $B$ and $j\in J_B$. Among additions of individual missing arc-product
coordinates, exactly $\rho_{Bj}$ are necessary and sufficient to determine
the full state circulation from the observed coordinates and linear balance
equations on the ambient circulation space. A minimum choice consists of
the nonforest edges of any spanning forest of $(V_B,U_{Bj})$.
The additional variables in \cref{thm:observed-compression} may instead be
chosen to represent these individual missing products, with the same
variable and row bounds and unit flow, product, and auxiliary coefficients.
\end{proposition}
\begin{proof}
The existing restriction map has nullity $\rho_{Bj}$. Each additional
scalar coordinate can raise its rank by at most one, proving necessity.
Let $F_{Bj}$ be a spanning forest of the unobserved subgraph, including
all its vertices. Add observations on $U_{Bj}\setminus F_{Bj}$, which
has $\rho_{Bj}$ edges. The remaining unobserved graph is a forest, so
the restriction kernel is zero by the proof of
\cref{thm:observed-compression}. This proves sufficiency and the stated
minimum.

Apply \cref{cor:forest-hull} to the completed observation set and regard
the added product coordinates as auxiliary variables. Projection onto the
original observations commutes with convexification, so this is an extended
formulation for the original $\Hull$. Completing labels only in $J_B$
does not introduce new labels. The added equations number
$\sum_{B,j}\rho_{Bj}\le\sum_Br_Ba_B$, already absorbed by the row
bound. The corollary's coefficient property includes the added coordinates.
\end{proof}

There is a direct interpretation of this completion. The block state flows
satisfy $A_Bf^{B,j}=y_jA_Bv_{E_B}$, since their deviations are block
circulations. Once the nonforest state products are known, sum these
equations on either side of a cut formed by deleting a remaining forest
edge from its component. That edge is the only
unknown crossing edge, so its state product is an affine expression in the
known state products and the state weight, with unit product coefficients.
Component balance equations enforce consistency if the unobserved forest
is disconnected. This is network elimination of the kind underlying reduced
RLT: \citet[Theorem~3.1 and Section~2]{LibertiPantelides2006} reconstruct
common-factor products by multiplying full-row-rank linear equations and
retaining products on suitable complementary coordinates. Here the extra
information is the observation-complement graph criterion and its exact
sparse simplex-hull consequence.

The minimum in \cref{prop:minimum-completion} concerns reconstruction by
\emph{individual coordinates on the ambient linear space}. It is neither
an extension-complexity lower bound nor a claim about the number of
coordinates needed on a capacity-degenerate slice. In particular,
$\lambda_j=0$ forces every state product to zero. The preprocessing in
\cref{lem:fixed-arcs} can remove fixed arcs before the parameters are
computed, but it does not remove this conditional-slice distinction.

All formulations also give explicit recovery. From retained coordinates,
recover each $t^{B,j}$ using \eqref{eq:pivot-reconstruction}, reconstruct
path deviations by $\widehat C_Bt^{B,j}$, and subtract their sum from
$s(x)$ to recover the merged state. Expand deviations along the original
paths using their signs. The proportional refinement
\eqref{eq:proportional-refinement} then supplies the global state flows
and the at-most-$m+1$-point decomposition from
\cref{prop:disaggregation}. This recovery is exact rational arithmetic for
rational input and a rational feasible extension; merely solving the
formulation numerically does not provide an exact certificate.

\begin{example}[Three observations on a $K_4$ block]\label{ex:forest-k4}
Orient every edge of $K_4$ from the smaller to the larger vertex and take
unit capacities and unit flow from vertex $1$ to vertex $4$. Thus
$b=(-1,0,0,1)^\top$. Use the reference flow $v_{14}=1$ and zero on the
other arcs. For one label $j$, observe $13$, $14$, and $24$; the unobserved
edges $12$, $23$, and $34$ form a spanning tree. Writing
$q_{13}=z_{13,j}$, $q_{14}=z_{14,j}-y_j$, and $q_{24}=z_{24,j}$, balance
gives the remaining state deviations
\[
 h_{12}=-q_{13}-q_{14},\qquad
 h_{23}=-q_{13}-q_{14}-q_{24},\qquad
 h_{34}=-q_{14}-q_{24}.
\]
For example, nonnegativity of the reconstructed flow on $23$ gives
\begin{equation}\label{eq:k4-forest-cut}
 z_{13,j}+z_{14,j}+z_{24,j}\le y_j.
\end{equation}
This cut is absent from independent product convexification. To see the
gap, take $m=1$, $y_1=1/2$, $z_{13,1}=z_{14,1}=z_{24,1}=1/5$, and let
$x$ average the four paths $1234$, $124$, $134$, and $14$ equally. Then
$x_{12}=x_{34}=1/2$ and the other four arc flows are $1/4$. All three
observed products satisfy \eqref{eq:mccormick}, but the left side of
\eqref{eq:k4-forest-cut} is $3/5>1/2$. The complete hull requires all
reconstructed bounds, including the residual bounds, not only this one cut.
\end{example}
